% RESEARCH_DRAFT_SHA256: 574ea551484cc409b704fc7db9a032c72bea91c7725f0f332958aafc3381529b
% Layout and source-reference wrappers expanded for the web renderer.
\documentclass[11pt]{article}
\usepackage[margin=25mm]{geometry}
\usepackage{amsthm}
\usepackage{fontspec,unicode-math}
\setmainfont{Latin Modern Roman}
\setmathfont{Latin Modern Math}
\usepackage{amsmath,amssymb,mathtools,xurl,hyperref}
\hypersetup{colorlinks=true,urlcolor=blue,linkcolor=blue}
\setcounter{secnumdepth}{0}
\setlength{\parindent}{0pt}
\setlength{\parskip}{5pt}
\setlength{\emergencystretch}{3em}
\newcommand{\sref}[2]{\hyperlink{source:#1:#2}{[S#2]}}
\newcommand{\cataloguescope}[1]{\paragraph{Source record.}#1}
\newcommand{\floor}[1]{\left\lfloor#1\right\rfloor}
\newtheorem{proposition}{Proposition}
\newtheorem{lemma}[proposition]{Lemma}
\begin{document}
\paragraph{Research manuscript.} OpenAI \texttt{gpt-6-astra}, reasoning \texttt{ultra}; 11 October 2026.
\paragraph{Preserved source body.} The complete catalogue mathematical text and sources follow with layout and source links expanded, including source qualifications and its historical preparation notes. New research begins at the section headed Research attempt.

% UnsolvedMath ID 20002203; source catalogue code AIM-INFRASTRUCTURE-0103
\section{A choice-point audit for social framing in mathematical practice}\label{20002203}
{\small UnsolvedMath ID 20002203 · Computer Science}
\subsection{Definitions and mathematical statement}

\paragraph{Shared notation.}Standard number systems are $\mathbb N,\mathbb Z,\mathbb Q,\mathbb R,\mathbb C$. Specialized notation and index ranges follow the cited source and the definitions below.


Mathematical creation includes choosing questions, representations, definitions, and arguments. A social influence question asks how people's experiences or assumptions affect those choices. A formal mathematical claim and an empirical account of how that claim came to be investigated are different objects of study.
\paragraph{Record scope and qualifications.}
The catalogue title is retained as metadata. Named certificates, contracts, bounds or frameworks in that title are not asserted as results of the original source extract.
The source names no corpus of mathematics, mechanism of influence, comparison design, or testable quantified hypothesis.

Catalogue screening limitation: Community, history, equity, or organizational policy discussion, rather than a mathematical statement.

The source snapshot labels this record partially solved; this is catalogue metadata rather than independent verification of a mathematical result or the current completion of a project.
\paragraph{Statement recorded in the supplied source.}
How Does the Gender Binary Influence the Math we Create?

How does the gender binary (and more generally how we perceive the world) influence the math we create? (not just how we teach and share it) \href{http://aimpl.org/gemscombin/7/}{[S2]}
\paragraph{Source status.}
The supplied catalogue labels this source record partially solved; that metadata was checked against the supplied snapshot on 2026-10-11. The label is preserved for provenance and is not an independent verification of present open status or mathematical correctness.
\subsection{Short English statement}
Investigate how perceptions of gender and the world may influence mathematical creation.
\subsection{Sources}
\begin{itemize}
\item[S1] Ulam AI and the UnsolvedMath Contributors, UnsolvedMath: A Curated Collection of Open Mathematics Problems, record 20002203 (AIM-INFRASTRUCTURE-0103). Catalogue curation and metadata: CC BY 4.0. The source extract is retained; definitions, English summaries and qualifications are editorial additions.
\url{https://huggingface.co/datasets/ulamai/UnsolvedMath}.
\item[S2] AIM Workshop Problem Lists, source linked by catalogue record AIM-INFRASTRUCTURE-0103.
\url{http://aimpl.org/gemscombin/7/}.
\end{itemize}
\paragraph{Preparation and review.}
The original source and UnsolvedMath contributors retain attribution for the source extract. Codex (GPT-6; exact serving revision not exposed) assisted with definitions, English summaries, scope notes and formatting. No new human mathematical authorship or independently verified solution is asserted. Permanent source ID 20002203 and code AIM-INFRASTRUCTURE-0103 checked against the complete supplied registry and current repository. Distinct upstream records, including repeated questions and broad or damaged extracts, are retained for complete coverage of the source CS category. A repeated record is not a new independent conjecture. Complete mathematical-equivalence and current-literature review has not been performed. Preparation and final review follow contribution prompts C and D, with source gaps disclosed. Structural and rendering checks do not establish a new theorem or open status.
\label{end:20002203}

\section{Research attempt}
\subsection{Literature, source context and scope}
The complete catalogue body above retains the original agenda, source status, and historical preparation. The AIMPL problem page could not be retrieved. The supplied complete source transcription was read, and the accessible primary AIM workshop announcement was checked for context. The announcement is not a substitute for verification of the inaccessible problem page. The catalogue's earlier proposed frameworks and research summaries are context, not established solutions credited to this attempt. The following is an explicitly limited investigation of one route; it neither substitutes a formalized toy model for the entire source agenda nor claims empirical or deployed-software results. Elementary auxiliary facts are not new original-target progress.

\subsection{An explicit information-loss route through a binary partition}
The original philosophical and empirical question asks how social framing influences mathematics. It does not specify a testable theorem about a given community. We examine a limited representational mechanism: retain only interactions between two categories and discard interactions within each. Fix four labeled vertices with partition $A=\{1,2\}$ and $B=\{3,4\}$. An encoder reports the partition sizes and the complete set of cross-partition edges, omitting same-partition edges.
\begin{proposition}
This encoding cannot determine even the maximum matching size of the original simple graph.
\end{proposition}
\begin{proof}
Let $G$ have no edges, and let $H$ have precisely edges $\{1,2\}$ and $\{3,4\}$. Both encodings report two vertices in each part and no cross-partition edges, so they are identical. In $G$ every matching has size zero. In $H$ the two declared edges share no endpoint, so they form a size-two matching; no matching on four vertices can have more than two edges, because each uses two distinct vertices. Thus the maximum matching sizes differ while the encoded data agree. Any function of the encoding alone must output the same value on both and hence fails for at least one graph.
\end{proof}
The failed inference is that a mathematical object built from only between-category relations retains every property of the original relations. The obstruction comes from the chosen loss of edges, not from a theorem that all two-category models are inappropriate: retaining all edges and the partition would avoid this particular loss. Nor does the example establish that the social gender binary caused any real mathematical choice.
\subsection{Verification and remaining gaps}
All edge subsets of the two tiny graphs were checked for disjoint endpoints, giving matching maxima 0 and 2 and identical encodings. No historical or sociological evidence linking social categories to mathematical creation was examined. The example supplies one logically possible representational consequence, with no causal explanation of the broad source question.

\section{References and source attribution}
Primary source: \url{http://aimpl.org/gemscombin/7/}. Attempted primary AIMPL page retrieval; failed with timeout or gateway error. The supplied full source transcription was read; the inaccessible live page was not independently inspected.

\url{https://aimath.org/pastworkshops/gemscombin.html}. Read the primary March 2023 workshop announcement and its community inclusion agenda.

All retrievals were checked on 11 October 2026. Original contributor and catalogue attribution remains in the preserved Sources section. The elementary arguments below are self-contained; their exposition is not a claim of first discovery.


\paragraph{Integration review.} Codex (GPT-6; exact serving revision not exposed) reviewed the written arguments and expanded layout and source-reference wrappers. Repository metadata and the production web renderer were checked using Python 3.13 and JavaScript with MathJax 3.2.2. This model-assisted review is not independent human proof verification; no PDF compilation or formal proof checking is asserted.

\section{Completion Estimate}
\noindent\textbf{COMPLETION: 0\%}
No substantiated new progress toward the entire original target is claimed. The worked results and failed route are delimited above; a known result, illustrative model, or auxiliary obstruction is not counted as original-target completion.

\section{Final status and attribution}
\textbf{Attempt status: unresolved.} This describes this research attempt, not a new ruling on the source catalogue's status or on external literature claims.
The mathematics and this self-audit were generated by OpenAI \texttt{gpt-6-astra}, reasoning \texttt{ultra}, on 11 October 2026. Codex (GPT-6; exact serving revision not exposed) orchestrated the batch and contributed formatting/repository review. No human mathematical authorship, independent human proof checking, or formal verification is claimed. The historical preparation paragraphs in the reproduced catalogue body concern that body only.

Contribution prompt A was applied to the complete definition and supplied source record, followed by the adversarial checks of prompt B. Local reproduction files are \texttt{batch013/checks.py}, \texttt{computations.json}, and the per-ID audit JSON. They were not public links at preparation time. Numerical checks are finite and do not prove unrestricted results. No LaTeX engine was found locally; compilation was not performed by this research agent.

\end{document}
